\begin{table}[htbp]
\centering
\caption{Standardized Effect Sizes}
\label{tab:sde}
\begin{adjustbox}{max width=\textwidth}
\begin{tabular}{lcccccc}
\hline\hline
Outcome & $\hat{\beta}$ & SE & SD($Y$) & SDE & SE(SDE) & Classification \\
\hline
\multicolumn{7}{l}{\textit{Panel A: Pooled}} \\[3pt]
All opioid deaths & -545.7 & 366.1 & 12.6 & -0.895 & 0.601 & Large negative \\
Rx opioid deaths & -35.1 & 71.1 & 2.8 & -0.263 & 0.534 & Large negative \\
Synthetic opioid deaths & -537.8 & 447.1 & 13.3 & -0.839 & 0.698 & Large negative \\
Rx opioid deaths (pop-weighted) & 58.8 & 71.5 & 2.8 & 0.442 & 0.537 & Large positive \\
[6pt]
\multicolumn{7}{l}{\textit{Panel B: Heterogeneous (Pre vs Post Fentanyl)}} \\[3pt]
All opioids (2015--2018) & 127.3 & 244.7 & 9.8 & 0.269 & 0.517 & Large positive \\
All opioids (2019--2022) & -635.0 & 307.5 & 13.8 & -0.951 & 0.461 & Large negative \\
\hline\hline
\multicolumn{7}{p{14cm}}{\footnotesize \textit{Notes:} \textbf{Country:} United States. \textbf{Research question:} Does disability prevalence causally increase opioid overdose mortality through insurance-mediated prescription access? \textbf{Policy mechanism:} SSDI/SSI disability enrollment provides Medicare (after 24-month waiting period) or Medicaid (immediately), granting prescription drug coverage that includes opioid analgesics, potentially creating a pipeline from disability adjudication to opioid prescribing. \textbf{Outcome definition:} Drug overdose deaths per 100,000 population by ICD-10 substance code from CDC VSRR provisional counts (12-month rolling totals). \textbf{Treatment:} Continuous; state-level disability prevalence rate (fraction of civilian noninstitutionalized population with a disability, ACS B18101). \textbf{Data:} CDC Vital Statistics Rapid Release (VSRR) drug overdose deaths by state-year-substance, 2015--2022; ACS 5-year estimates; 41 states, 261 state-year observations. \textbf{Method:} OLS with state and year fixed effects, standard errors clustered at state level. Within-state panel exploiting temporal variation in disability prevalence. \textbf{Sample:} 41 states with non-suppressed mortality counts across all years; excludes territories and states with systematic CDC data suppression. SDE $= \hat{\beta} \times \text{SD}(X) / \text{SD}(Y)$ for continuous treatment. Classification refers to magnitude, not statistical significance: Large ($|$SDE$| > 0.15$), Moderate ($0.05$--$0.15$), Small ($0.005$--$0.05$), Null ($< 0.005$).} \\
\end{tabular}
\end{adjustbox}
\end{table}
